\section{An exponential rank-loss bound}\label{sec:rank}

We now bound the rank lost by an idempotent diagram whose image adds
one off-diagonal pair to the identity relation. The amplification step
uses $32$ conjugates of that diagram to construct $256$ successive
relation additions. Their joint support bounds the total rank loss,
while a smaller relation corner bounds the loss at each addition.

For every integer $h\ge2$, set
\begin{equation}\label{eq:rank-loss-function}
 \Lfun(h)=2^{\lfloor(h-2)/31\rfloor}.
\end{equation}
All reductions below are the maps $R_e$ of Lemma~\ref{lem:corners}.
Their ranges are regarded as submonoids of the indicated smaller
matching monoids. In particular, their identities are the ambient
identity diagrams, and they preserve absolute rank.

The theorem first assumes permutation lifts in addition to a unital
surjection onto the \emph{entire} relation monoid. A surjection onto only
the relations generated by the displayed additions would not suffice:
the proof must also lift arbitrary elements of smaller relation corners.
Corollary~\ref{cor:divisor-degree} removes the extra lift assumption by
passing to a minimum-rank idempotent corner.

\begin{theorem}[Rank loss for a single relation addition]
\label{thm:rank-loss}
Let $H$ be a finite set of size $h\ge2$, let $m\ge0$, and let
$S\subseteq\D{m}$ be a submonoid containing $\id$.
Suppose that $\phi:S\to\R{H}$ is a unital surjection and that every
permutation of $H$ has a lift in $S$ which is a permutation diagram.
If $a\in S$ is idempotent and
\[
 \phi(a)=\Id{H}\cup\{(x,y)\},\qquad x,y\in H,\quad x\ne y,
\]
then
\[
 m-\rank(a)\ge\Lfun(h).
\]
\end{theorem}

\begin{proof}
We induct on $h$, simultaneously for every ambient degree $m$ and
every choice of $S,\phi,a$ satisfying the hypotheses. Put $t=16$.
For $2\le h<2t+1=33$, one has $\Lfun(h)=1$. A full-rank diagram is
a permutation, and an idempotent permutation is the identity.
Since $\phi$ is unital and $\phi(a)\ne\Id{H}$, the rank of $a$ is
strictly smaller than $m$. This proves the base cases.

Now let $h\ge33$ and assume the theorem for every smaller size at
least two. Write
\[
 c=m-\rank(a).
\]

\paragraph{Conjugates with a common support bound.}
Choose pairwise disjoint sets
\[
 U=\{u_1,\ldots,u_t\},\qquad
 V=\{v_1,\ldots,v_t\},\qquad \{z\}
\]
in $H$. For each ordered pair of distinct points of $H$, a
permutation $\pi$ can send $(x,y)$ to that pair. Choose a
permutation-diagram lift $\widetilde\pi\in S$. Its inverse belongs to
$S$ because a permutation diagram has finite order. In path order,
\[
 \phi(\widetilde\pi^{-1}a\widetilde\pi)
 =\pi^{-1}\bigl(\Id{H}\cup\{(x,y)\}\bigr)\pi
 =\Id{H}\cup\{(\pi(x),\pi(y))\}.
\]
Using the pairs $(u_i,z)$ and $(z,v_j)$ in this construction gives
idempotents $s_i,s'_j\in S$
such that
\begin{equation}\label{eq:rank-conjugate-images}
 \phi(s_i)=\Id{H}\cup\{(u_i,z)\},\qquad
 \phi(s'_j)=\Id{H}\cup\{(z,v_j)\}
 \quad(1\le i,j\le t).
\end{equation}
Conjugation by permutation diagrams preserves rank, so all these
idempotents have rank $m-c$. By Lemma~\ref{lem:corners}(ii), each
has support of size at most $2c$. Thus the set
\[
 J=\bigcup_{i=1}^{t}\supp(s_i)
       \ \cup\ \bigcup_{j=1}^{t}\supp(s'_j)
 \subseteq[m]
\]
has size at most $4tc$. Every product $s_i s'_j$ is supported in $J$,
because diagrams supported in a fixed set form a submonoid.

\paragraph{A chain of $t^2$ additions.}
Put $P_0=\Id{H\setminus\{z\}}$. Choose an idempotent $b_0\in S$
with $\phi(b_0)=P_0$: take any lift of $P_0$ and then a positive
idempotent power. The image stays $P_0$ because $P_0$ is idempotent.
For $1\le i,j\le t$, define
\[
 g_{ij}=b_0s_i s'_j b_0\in b_0Sb_0.
\]
In path order, the two added pairs in
\eqref{eq:rank-conjugate-images} allow the path $u_i,z,v_j$.
The outer factors $P_0$ delete pairs with endpoint $z$, so
\begin{equation}\label{eq:rank-generator-image}
 \phi(g_{ij})=P_0\cup\{(u_i,v_j)\}.
\end{equation}
Figure~\ref{fig:amplification} shows how the same $2t$ conjugates provide
all $t^2$ choices of the pair $(u_i,v_j)$.

\begin{figure}[htbp]
\centering
\begin{tikzpicture}[>=Stealth,vertex/.style={circle,draw,minimum size=7mm,inner sep=1pt},every node/.style={font=\small}]
\node[vertex] (upperleft) at (0,1.2) {$u_1$};
\node[vertex] (middleleft) at (0,0) {$u_i$};
\node[vertex] (lowerleft) at (0,-1.2) {$u_t$};
\node at (0,0.6) {$\vdots$};
\node at (0,-0.6) {$\vdots$};
\node[vertex] (hub) at (2.2,0) {$z$};
\node[vertex] (upperright) at (4.4,1.2) {$v_1$};
\node[vertex] (middleright) at (4.4,0) {$v_j$};
\node[vertex] (lowerright) at (4.4,-1.2) {$v_t$};
\node at (4.4,0.6) {$\vdots$};
\node at (4.4,-0.6) {$\vdots$};
\draw[->,gray] (upperleft)--(hub);
\draw[->,very thick,blue!65!black] (middleleft)--node[above] {$\phi(s_i)$} (hub);
\draw[->,gray] (lowerleft)--(hub);
\draw[->,gray] (hub)--(upperright);
\draw[->,very thick,blue!65!black] (hub)--node[above] {$\phi(s'_j)$} (middleright);
\draw[->,gray] (hub)--(lowerright);
\node[align=center] at (1.05,1.7) {incoming\\family};
\node[align=center] at (3.35,1.7) {outgoing\\family};
\draw[->,thick] (5.15,0)--node[above,align=center] {sandwich\\by $P_0$} (6.65,0);
\node[vertex] (retainedleft) at (7.35,0) {$u_i$};
\node[vertex] (retainedright) at (9.2,0) {$v_j$};
\draw[->,very thick,blue!65!black] (retainedleft)--(retainedright);
\node[align=center] at (8.25,-0.9) {$\phi(g_{ij})=P_0\cup\{(u_i,v_j)\}$};
\end{tikzpicture}
\caption{The gray arrows are alternatives from the two families of
relation images; each individual image has one added arrow. The
highlighted path creates the selected pair after both outer factors
delete endpoints at $z$. Identity loops are omitted. Only relations on
$H$ are shown. The common matching support bound comes from taking the
union of the supports of the $32$ conjugates and reducing at $b_0$; the same
conjugates provide all $256$ choices.}
\label{fig:amplification}
\end{figure}


Let $r_0=\rank(b_0)$. Applying Lemma~\ref{lem:corners}(iii) at
$b_0$ to the single set $J$ gives one set $J_0\subseteq[r_0]$,
with $|J_0|\le4tc$, such that
\begin{equation}\label{eq:rank-common-support}
 \supp\bigl(R_{b_0}(g_{ij})\bigr)\subseteq J_0
 \qquad(1\le i,j\le t).
\end{equation}
This set is common to all $t^2$ generators.

List the pairs $(i,j)\in[t]\times[t]$ in any order, without
repetition. Let $E_\ell$ be the set of relation pairs $(u_i,v_j)$
in the first $\ell$ positions, and set
\[
 P_\ell=P_0\cup E_\ell\qquad(0\le\ell\le t^2).
\]
All additional pairs go from $U$ to $V$, and $U\cap V=\varnothing$.
Consequently two additional pairs cannot be composed with each other;
the identity pairs in $P_0$ preserve their endpoints. Thus
\begin{equation}\label{eq:rank-union}
 (P_0\cup A)(P_0\cup B)=P_0\cup(A\cup B)
 \qquad(A,B\subseteq U\times V).
\end{equation}
In particular, $P_\ell^2=P_\ell$.

Starting from $b_0$, construct $b_\ell$ as follows. If $(i,j)$ occupies
position $\ell$, take $b_\ell$ to be a positive idempotent power of
\[
 b_{\ell-1}g_{ij}b_{\ell-1}.
\]
Inductively, its image before taking the power is $P_\ell$, by
\eqref{eq:rank-generator-image} and \eqref{eq:rank-union};
the power has the same image. Hence
\begin{equation}\label{eq:rank-chain}
 b_\ell^2=b_\ell,\qquad
 b_\ell\in b_{\ell-1}Sb_{\ell-1},\qquad
 \phi(b_\ell)=P_\ell.
\end{equation}
Every $b_\ell$ also belongs to $b_0Sb_0$. Indeed, this is true
initially, and this latter set is a submonoid containing every $g_{ij}$;
it is therefore preserved by the construction.

Write
\[
 r_\ell=\rank(b_\ell)\quad(0\le\ell\le t^2),
 \qquad
 \Delta_\ell=r_{\ell-1}-r_\ell\quad(1\le\ell\le t^2).
\]
The nested corners in \eqref{eq:rank-chain} and
Lemma~\ref{lem:rank-basic} give $\Delta_\ell\ge0$.
In the single coordinate system obtained
by reducing at $b_0$, the starting element is the identity and all
generators are supported in $J_0$ by \eqref{eq:rank-common-support}.
The same is true of every constructed product and power, in particular
$R_{b_0}(b_{t^2})$. Rank preservation and the elementary bound of rank
defect by support size give
\begin{equation}\label{eq:rank-total-budget}
 \sum_{\ell=1}^{t^2}\Delta_\ell
 =r_0-r_{t^2}
 \le\bigl|\supp(R_{b_0}(b_{t^2}))\bigr|
 \le|J_0|\le4tc.
\end{equation}

\paragraph{The loss at one addition.}
We claim that for every $1\le\ell\le t^2$,
\begin{equation}\label{eq:rank-single-step}
 2\Delta_\ell\ge\Lfun(h-2t+1).
\end{equation}
Each local argument uses fresh port coordinates to certify this bound
on the already defined absolute rank difference $\Delta_\ell$ in
\eqref{eq:rank-total-budget}.
Fix $\ell$, let $(i,j)$ be its pair, and put
$P=P_{\ell-1}$. Reduce at $b_{\ell-1}$ and define
\[
 T=R_{b_{\ell-1}}(b_{\ell-1}Sb_{\ell-1})
       \subseteq\D{r_{\ell-1}},\qquad
 \beta=R_{b_{\ell-1}}(b_\ell)\in T.
\]
Here $T$ contains the ambient identity
$R_{b_{\ell-1}}(b_{\ell-1})=\id$. The unital-surjection assertion of
Lemma~\ref{lem:unit-lifts} gives
\[
 \theta:T\longrightarrow P\R{H}P,
 \qquad
 \theta(R_{b_{\ell-1}}(w))=\phi(w)
 \quad(w\in b_{\ell-1}Sb_{\ell-1}).
\]
Moreover, $\beta$ is idempotent,
$\theta(\beta)=P_\ell$, and rank preservation gives
\begin{equation}\label{eq:rank-step-support}
 \rank(\beta)=r_\ell,
 \qquad |\supp(\beta)|\le2\Delta_\ell,
\end{equation}
where support now refers to the indices $[r_{\ell-1}]$.

\paragraph{A full relation corner on fewer points.}
Set
\[
 F=\{u_i,v_j\}\cup\bigl(H\setminus(U\cup V\cup\{z\})\bigr),
 \qquad E=\Id{F},\qquad Q=PEP.
\]
Thus $|F|=h-2t+1=h-31$. Among the possible extra pairs in $P$, the
only one with both endpoints in $F$ is $(u_i,v_j)$, which has not yet
been added. Every point of $F$ is different from $z$. Therefore
\[
 EPE=E.
\]
Lemma~\ref{lem:relation-corner} shows that $Q$ is idempotent,
$PQ=QP=Q$, and restriction $Z\mapsto EZE$ is a unital isomorphism
from $Q\R{H}Q$ onto $\R{F}$. In particular,
\[
 Q(P\R{H}P)Q=Q\R{H}Q.
\]
Equation~\eqref{eq:rank-union} also gives $PP_\ell=P_\ell P=P_\ell$.
Since $EQ=EP$ and $QE=PE$, we obtain
\begin{equation}\label{eq:rank-smaller-image}
 E(QP_\ell Q)E
 =EP P_\ell PE
 =EP_\ell E
 =\Id{F}\cup\{(u_i,v_j)\}.
\end{equation}

\paragraph{The second reduction and the induction hypotheses.}
Among the idempotents $d\in T$ with $\theta(d)=Q$, choose one of
minimum rank. Such idempotents exist because $\theta$ is onto and
an idempotent power of a lift of $Q$ still maps to $Q$. This minimum
is taken inside $T$, in its current coordinates in
$\D{r_{\ell-1}}$. Put $r'=\rank(d)$ and reduce once more, defining
\[
 T'=R_d(dTd)\subseteq\D{r'}.
\]
The identity of $T'$ is $R_d(d)=\id$. Its induced map onto
$Q(P\R{H}P)Q$, followed by restriction to $F$, is the unital
surjection
\[
 \psi:T'\longrightarrow\R{F},\qquad
 \psi(R_d(w))=E\theta(w)E\quad(w\in dTd).
\]
By Lemma~\ref{lem:unit-lifts} and the minimum-rank choice of $d$,
every unit of $Q(P\R{H}P)Q$ has a permutation-diagram lift in $T'$.
The restriction isomorphism preserves units, so every permutation of
$F$ has such a lift under $\psi$.

Let $J_\ell=\supp(\beta)\subseteq[r_{\ell-1}]$. Apply
Lemma~\ref{lem:corners}(iii) to the idempotent $d$ in this ambient
matching monoid. There is a set $K\subseteq[r']$ with
$|K|\le|J_\ell|\le2\Delta_\ell$ such that
\[
 \supp\bigl(R_d(d\beta d)\bigr)\subseteq K.
\]
The sandwich $d\beta d$ need not be idempotent. Choose a positive
idempotent power $v$ of its reduction $R_d(d\beta d)$ in $T'$.
Support remains in $K$ under taking powers. By
\eqref{eq:rank-smaller-image}, the image before taking the power is
$\Id{F}\cup\{(u_i,v_j)\}$, an idempotent relation; hence
\[
 v^2=v,
 \qquad \psi(v)=\Id{F}\cup\{(u_i,v_j)\},
 \qquad |\supp(v)|\le2\Delta_\ell.
\]
We have obtained an identity-containing submonoid of $\D{r'}$, a
unital surjection onto $\R{F}$ with permutation lifts, and the
required idempotent above a single off-diagonal addition. These are
all the hypotheses of the theorem at size $|F|=h-31$, which lies
between $2$ and $h-1$. The induction hypothesis and the support bound
therefore give
\[
 \Lfun(h-31)\le r'-\rank(v)
        \le|\supp(v)|\le2\Delta_\ell.
\]
This proves \eqref{eq:rank-single-step} for the chosen $\ell$, and
hence for every step of the chain.

\paragraph{Completion of the recurrence.}
The bounds established above are collected in
Table~\ref{tab:rank-budget}. The support is measured once, in the
coordinates of $b_0$; the separate smaller corners certify each
summand in the resulting rank-loss budget.

\begin{table}[htbp]
\centering
\renewcommand{\arraystretch}{1.22}
\begin{tabular}{lcc}
\hline
Quantity & General expression & $t=16$ \\
\hline
Conjugates & $2t$ & $32$ \\
Successive additions & $t^2$ & $256$ \\
Common support budget & $4tc$ & $64c$ \\
Retained relation points & $h-2t+1$ & $h-31$ \\
Loss per addition & $\frac12\Lfun(h-2t+1)$ & $\frac12\Lfun(h-31)$ \\
Total required loss & $\frac{t^2}{2}\Lfun(h-2t+1)$ & $128\Lfun(h-31)$ \\
\hline
\end{tabular}
\caption{The support upper bound and the inductive rank-loss lower
bound are applied to the same chain. Here $c=m-\rank(a)$.}
\label{tab:rank-budget}
\end{table}

Summing \eqref{eq:rank-single-step} and applying
\eqref{eq:rank-total-budget} yields
\[
 4tc\ge\sum_{\ell=1}^{t^2}\Delta_\ell
      \ge\frac{t^2}{2}\Lfun(h-2t+1).
\]
With $t=16$, it follows that
\[
 c\ge2\Lfun(h-31)
   =2^{1+\lfloor(h-33)/31\rfloor}
   =2^{\lfloor(h-2)/31\rfloor}
   =\Lfun(h).
\]
This completes the induction.
\end{proof}

\begin{corollary}[Degree bound for a unital relation-monoid divisor]
\label{cor:divisor-degree}
Let $H$ be a finite set of size $h\ge2$. If an identity-containing
submonoid $S\subseteq\D{k}$ admits a unital surjection
$\phi:S\to\R{H}$, then
\[
 k\ge\Lfun(h)=2^{\lfloor(h-2)/31\rfloor}.
\]
No hypothesis about permutation lifts is required here.
\end{corollary}

\begin{proof}
Choose an idempotent $e\in S$ of minimum rank among those with
$\phi(e)=\Id{H}$. The set of choices is nonempty, since the identity
of $S$ is one such idempotent. Put $r=\rank(e)$.
By Lemmas~\ref{lem:corners} and~\ref{lem:unit-lifts}, the reduced
corner
\[
 T=R_e(eSe)\subseteq\D{r}
\]
contains the ambient identity and maps unitally onto
$\Id{H}\R{H}\Id{H}=\R{H}$. The minimum-rank choice supplies
permutation-diagram lifts for all permutations of $H$.

Choose distinct $x,y\in H$. Surjectivity supplies a lift in $T$ of
$\Id{H}\cup\{(x,y)\}$. A positive idempotent power of that lift is
an idempotent $a\in T$ with the same image. Theorem~\ref{thm:rank-loss}
applies to this representation, so
\[
 \Lfun(h)\le r-\rank(a)\le r\le k.\qedhere
\]
\end{proof}
